\section{Conclusion}

TRACE-C is best understood as a strictly-prior rank-calibrated detector, not a
literal copula calibration procedure. It combines magnitude-preserving robust
residuals, local/dependence/temporal channels, a rank-aggregated score, and an
explicit selection path. Its rank $p$-values support ordering and selection;
they are not estimated event probabilities or posterior confidence.

The evidence is mixed and operationally informative. TRACE-C ranks Storm
Atiyah first in 2019 development, but the channel ablation shows that this is
a local-led window: copula-only ranks it 59th, and dropping the copula channel
leaves it 2nd without a record alert. Simple reconstruction, and the temporal
channel alone, rank the short frequency event much better than the fused
detector. The frozen 2020 run selects no alerts, while its ranked list contains
plausible weather-associated and unlabelled extremes. The result exposes record
saturation, fusion trade-offs, and sensor/aggregation limits rather than
establishing nominal FDR or copula discovery.

Future work should pre-register a dependence-aware selector, evaluate it on a
new untouched period, preserve sub-period frequency information, test missing
streams and alternative aggregations, and replace fixed 48-row attention blocks
with calendar-aware budgeting. The current contribution is the narrower one:
an auditable algorithm and evidence package whose assumptions, fallbacks, and
negative results are stated at the same resolution as its successes.
